\section{Results}

\subsection{Machine Translation}

The Transformer improves over strong prior baselines on both WMT 2014 English-to-German and English-to-French. The larger model achieves a new state of the art on English-to-German and remains highly competitive on English-to-French while requiring substantially less training cost than earlier recurrent or convolutional systems.

\begin{table}[h]
\centering
\caption{Headline comparison on WMT 2014 translation benchmarks.}
\label{tab:headline-results}
\begin{tabular}{lcc}
\toprule
Model & EN-DE BLEU & EN-FR BLEU \\
\midrule
GNMT + RL & 24.6 & 39.92 \\
ConvS2S & 25.16 & 40.46 \\
Transformer (base) & 27.3 & 38.1 \\
Transformer (big) & 28.4 & 41.0 \\
\bottomrule
\end{tabular}
\end{table}

These results indicate that the attention-centric formulation is effective in practice. In particular, the model appears to capture long-range dependencies while remaining easy to optimize in parallel.

\subsection{Model Variations}

We also tested several variations of the Transformer family during development. The overall pattern was consistent with the headline results: larger attention-based models performed better, and the architecture was robust enough to support a strong translation system across both benchmark settings.

\subsection{Interpretability}

Qualitatively, some attention heads appear to focus on linguistically meaningful relationships. This suggests that the model may also offer a more interpretable view of sequence transduction behavior than purely recurrent alternatives.
